%Página
\section{Prático}

\begin{enumerate}
	\item CA Root / CA Intermediária - Estrutura de Diretórios
	\begin{tcolorbox}[width=\textwidth]
	{\footnotesize
		\begin{verbatim}
			$ mkdir -p /root/ca/root-ca/{certs,crl,newcerts,private}
			$ mkdir -p /root/ca/intermediate-ca/{certs,crl,csr,newcerts,private}
			$ mkdir /root/ca/user-certs
			$ chmod 700 /root/ca/root-ca/private /root/ca/intermediate-ca/private
			$ touch /root/ca/root-ca/index.txt /root/ca/intermediate-ca/index.txt
			$ echo 1000 > /root/ca/root-ca/serial
			$ echo 1000 > /root/ca/intermediate-ca/serial
			$ echo 1000 > /root/ca/root-ca/crlnumber
			$ echo 1000 > /root/ca/intermediate-ca/crlnumber
		\end{verbatim}
	}
	\end{tcolorbox}

	\item CA Root - Configuração da CA Root (/root/ca/root-ca/openssl.cnf)
	\begin{tcolorbox}[width=\textwidth, enhanced, breakable]
	{\footnotesize
		\begin{verbatim}
			[ ca ]
			default_ca = CA_default
			
			[ CA_default ]
			dir               = /root/ca/root-ca
			certs             = $dir/certs
			crl_dir           = $dir/crl
			new_certs_dir     = $dir/newcerts
			database          = $dir/index.txt
			serial            = $dir/serial
			private_key       = $dir/private/root-ca.key.pem
			certificate       = $dir/certs/root-ca.cert.pem
			crlnumber         = $dir/crlnumber
			crl               = $dir/crl/root-ca.crl.pem
			crl_extensions    = crl_ext
			default_crl_days  = 30
			default_md        = sha256
			name_opt          = ca_default
			cert_opt          = ca_default
			default_days      = 3650
			preserve          = no
			policy            = policy_strict
			
			[ policy_strict ]
			countryName             = match
			stateOrProvinceName     = match
			organizationName        = match
			organizationalUnitName  = optional
			commonName              = supplied
			emailAddress            = optional
			
			[ req ]
			default_bits        = 4096
			distinguished_name  = req_distinguished_name
			string_mask         = utf8only
			default_md          = sha256
			prompt              = no
			
			[ req_distinguished_name ]
			countryName                     = BR
			stateOrProvinceName             = Sao Paulo
			localityName                    = Sao Paulo
			0.organizationName              = senhasegura
			organizationalUnitName          = IT Department
			commonName                      = My Root CA
			emailAddress                    = mt4@senhasegura.com
			
			[ v3_ca ]
			subjectKeyIdentifier = hash
			authorityKeyIdentifier = keyid:always,issuer
			basicConstraints = critical,CA:true
			keyUsage = critical, digitalSignature, cRLSign, keyCertSign
			
			[ v3_intermediate_ca ]
			subjectKeyIdentifier = hash
			authorityKeyIdentifier = keyid:always,issuer
			basicConstraints = critical,CA:true,pathlen:0
			keyUsage = critical, digitalSignature, cRLSign, keyCertSign
			
			[ crl_ext ]
			authorityKeyIdentifier=keyid:always
		\end{verbatim}
	}
	\end{tcolorbox}

	\item CA Intermediária - Configuração da CA Intermediária (/root/ca/intermediate-ca/openssl.cnf)
	\begin{tcolorbox}[width=\textwidth, enhanced, breakable]
	{\footnotesize
		\begin{verbatim}
			$ cat /root/ca/intermediate-ca/openssl.cnf

			[ ca ]
			default_ca = CA_default
			
			[ CA_default ]
			dir               = /root/ca/intermediate-ca
			certs             = $dir/certs
			crl_dir           = $dir/crl
			new_certs_dir     = $dir/newcerts
			database          = $dir/index.txt
			serial            = $dir/serial
			private_key       = $dir/private/intermediate-ca.key.pem
			certificate       = $dir/certs/intermediate-ca.cert.pem
			crlnumber         = $dir/crlnumber
			crl               = $dir/crl/intermediate-ca.crl.pem
			crl_extensions    = crl_ext
			default_crl_days  = 30
			default_md        = sha256
			name_opt          = ca_default
			cert_opt          = ca_default
			default_days      = 3650
			preserve          = no
			policy            = policy_loose
			
			[ policy_loose ]
			countryName             = match
			stateOrProvinceName     = match
			organizationName        = match
			organizationalUnitName  = optional
			commonName              = supplied
			emailAddress            = optional
			
			[ req ]
			default_bits        = 4096
			distinguished_name  = req_distinguished_name
			string_mask         = utf8only
			default_md          = sha256
			prompt              = no
			
			[ req_distinguished_name ]
			countryName                     = BR
			stateOrProvinceName             = Sao Paulo
			localityName                    = Sao Paulo
			0.organizationName              = senhasegura
			organizationalUnitName          = SEGi9
			commonName                      = My Intermediate CA
			emailAddress                    = segi9@senhasegura.com
			
			[ v3_intermediate_ca ]
			subjectKeyIdentifier = hash
			authorityKeyIdentifier = keyid:always,issuer
			basicConstraints = critical,CA:true,pathlen:0
			keyUsage = critical, digitalSignature, cRLSign, keyCertSign
			
			[ crl_ext ]
			authorityKeyIdentifier=keyid:always
			
			[ usr_cert ]
			basicConstraints = CA:FALSE
			nsCertType = client, email
			nsComment = "User Certificate"
			subjectKeyIdentifier = hash
			authorityKeyIdentifier = keyid,issuer
			keyUsage = critical, nonRepudiation, digitalSignature, keyEncipherment
			extendedKeyUsage = clientAuth, emailProtection
		\end{verbatim}
	}
	\end{tcolorbox}

	\item CA Root - Criar o Par de Chaves e Certificado da CA Root
	\begin{tcolorbox}[width=\textwidth]
	{\footnotesize
		\begin{verbatim}
			$ openssl genpkey -algorithm RSA -out /root/ca/root-ca/private/root-ca.key.pem \
			-pkeyopt rsa_keygen_bits:4096

			$ chmod 400 /root/ca/root-ca/private/root-ca.key.pem
			
			$ openssl req -config /root/ca/root-ca/openssl.cnf -key /root/ca/root-ca/private/root-ca.key.pem \
			-new -x509 -days 7300 -sha256 -extensions v3_ca -out /root/ca/root-ca/certs/root-ca.cert.pem

			$ chmod 444 /root/ca/root-ca/certs/root-ca.cert.pem
		\end{verbatim}
	}
	\end{tcolorbox}

	\item CA Intermediária - Criar o Par de Chaves da CA Intermediária
	\begin{tcolorbox}[width=\textwidth]
	{\footnotesize
		\begin{verbatim}
			$ openssl genpkey -algorithm RSA -out /root/ca/intermediate-ca/private/intermediate-ca.key.pem \
			-pkeyopt rsa_keygen_bits:4096

			$ chmod 400 /root/ca/intermediate-ca/private/intermediate-ca.key.pem
		\end{verbatim}
	}
	\end{tcolorbox}

	\item CA Intermediária - Criar a Solicitação de Certificado (CSR) da CA Intermediária
	\begin{tcolorbox}[width=\textwidth]
	{\footnotesize
		\begin{verbatim}
			$ openssl req -config /root/ca/intermediate-ca/openssl.cnf -new -sha256 \
			-key /root/ca/intermediate-ca/private/intermediate-ca.key.pem \
			-out /root/ca/intermediate-ca/csr/intermediate-ca.csr.pem
		\end{verbatim}
	}
	\end{tcolorbox}

	\item CA Root / CA Intermediária - Assinar o Certificado da CA Intermediária com a CA Root
	\begin{tcolorbox}[width=\textwidth]
	{\footnotesize
		\begin{verbatim}
			$ openssl ca -batch -config /root/ca/root-ca/openssl.cnf -extensions v3_intermediate_ca \
			-days 3650 -notext -md sha256 -in /root/ca/intermediate-ca/csr/intermediate-ca.csr.pem \
			-out /root/ca/intermediate-ca/certs/intermediate-ca.cert.pem

			$ chmod 444 /root/ca/intermediate-ca/certs/intermediate-ca.cert.pem
		\end{verbatim}
	}
	\end{tcolorbox}

	\item CA Intermediária - Criar a Cadeia de Certificados
	\begin{tcolorbox}[width=\textwidth]
	{\footnotesize
		\begin{verbatim}
			$ cat /root/ca/intermediate-ca/certs/intermediate-ca.cert.pem /root/ca/root-ca/certs/root-ca.cert.pem > \
			/root/ca/intermediate-ca/certs/ca-chain.cert.pem

			$ chmod 444 /root/ca/intermediate-ca/certs/ca-chain.cert.pem
		\end{verbatim}
	}
	\end{tcolorbox}

	\item Usuário - Criar o Par de Chaves do Usuário Final
	\begin{tcolorbox}[width=\textwidth]
	{\footnotesize
		\begin{verbatim}
			$ openssl genpkey -algorithm RSA -out /root/ca/user-certs/user.key.pem -pkeyopt rsa_keygen_bits:2048

			$ chmod 400 /root/ca/user-certs/user.key.pem
		\end{verbatim}
	}
	\end{tcolorbox}

	\item Usuário - Criar a Solicitação de Certificado (CSR) do Usuário Final
	\begin{tcolorbox}[width=\textwidth]
	{\footnotesize
		\begin{verbatim}
			$ openssl req -new -sha256 -key /root/ca/user-certs/user.key.pem \
			-subj "/C=BR/ST=Sao Paulo/L=Sao Paulo/O=senhasegura/OU=Users/CN=usuario.exemplo" \
			-out /root/ca/user-certs/user.csr.pem
		\end{verbatim}
	}
	\end{tcolorbox}

	\item Usuário - Assinar o Certificado de Usuário Final com a CA Intermediária
	\begin{tcolorbox}[width=\textwidth]
	{\footnotesize
		\begin{verbatim}
			$ openssl ca -batch -config /root/ca/intermediate-ca/openssl.cnf -extensions usr_cert -days 375 -notext \
			-md sha256 -in /root/ca/user-certs/user.csr.pem -out /root/ca/user-certs/user.cert.pem

			$ chmod 444 /root/ca/user-certs/user.cert.pem
		\end{verbatim}
	}
	\end{tcolorbox}

	\item Usuário - Verificar a Cadeia de Certificados
	\begin{tcolorbox}[width=\textwidth]
	{\footnotesize
		\begin{verbatim}
			$ openssl verify -CAfile /root/ca/intermediate-ca/certs/ca-chain.cert.pem \
			/root/ca/user-certs/user.cert.pem
		\end{verbatim}
	}
	\end{tcolorbox}

\end{enumerate}

\pagebreak

%Copyright
\vspace*{\fill}
\begin{flushright}
	\underline{\textit{Copyright}}\\
	All rights reserved to MT4 Tecnologia LTDA.\bigskip

	\underline{\textit{Propósito deste documento}}\\
	Infraestrutura de Chaves Públicas (PKI)\bigskip
	
	\underline{\textit{MT4 Tecnologia Ltda.}}\\
	Rua Joaquim Antunes, 767 – Pinheiros – São Paulo\\
	Phone: +55 11 3065-7070\\
	\url{www.mt4.com.br}\\
	\href{mailto:support@senhasegura.com}{support@senhasegura.com}\\
\end{flushright}
